% !TEX program = pdflatex
% !TEX enableSynctex = true
% !BIB program = bibtex

\documentclass[12pt]{article}

\usepackage{setspace}
\usepackage{amsmath}
\usepackage{amsfonts}
\usepackage{graphicx}
\usepackage{float}
\usepackage{natbib}
\addtolength{\oddsidemargin}{-.7in}
\addtolength{\evensidemargin}{-.7in}
\addtolength{\textwidth}{1.4in}
\usepackage{enumerate}
\onehalfspacing
\usepackage{geometry} % Required for customizing page layout
\usepackage{ragged2e}

\usepackage{caption}
\usepackage{booktabs}

\usepackage{hyperref}
\hypersetup{
	pdfstartview = FitH,
	pdfauthor = {...},
	pdftitle = {...},
	pdfkeywords = {...; ...; ...; ...},
	colorlinks = true,
	linkcolor = blue,
	urlcolor = blue,
	citecolor = blue,
	linktocpage=true
}

\title{Bankruptcy Resolution and the Macroeconomics of Cash Flow Based Lending}
\author{Barnabás Székely \\ \small  Goethe-Universität Frankfurt}
\date{}

\begin{document}

\maketitle


\section{Introduction}

This paper argues that high reorganization costs limit access to cash flow-backed borrowing, distort firm financing, and depress aggregate productivity. I develop a general equilibrium model with heterogeneous firms that may borrow against physical assets or expected future cash flows and face bankruptcy outcomes reflecting Chapter 7 (liquidation) and Chapter 11 (reorganization) of the U.S. Bankruptcy Code. Consistent with the empirical evidence, large administrative and legal costs of reorganization often force small and medium-sized firms to liquidate, whereas large firms can typically reorganize in default. Since liquidation wipes out all future cash flows, elevated liquidation risk prevents smaller firms from credibly pledging future earnings as collateral. As a result, asset-poor firms are simultaneously constrained by insufficient physical collateral and high liquidation risk. This `financing trap' distorts firm financing, inhibits entry and growth, which can potentially translate into significant aggregate productivity losses.

To quantify the broader macroeconomic implications of this mechanism, I develop a general equilibrium model that captures the effects of reorganization costs on firms' liquidation risk, access to cash flow-backed debt, and the resulting consequences for capital allocation and aggregate productivity. The model features heterogeneous firms, in-equilibrium defaults, and firm-specific external finance premia (Khan, Senga, and Thomas, 2017), and integrates endogenous bankruptcy resolution decisions (Tamayo, 2017; Antill and Grenadier, 2019; Corbae and D'Erasmo, 2021) with heterogeneity in debt contracts (Lian and Ma, 2021; Drechsel, 2023). To the best of my knowledge, it is the first model to analyze the macroeconomic effects of credit frictions in an environment where firms can access both asset-based and cash flow-based debt while facing endogenous, size-dependent liquidation risk.

\section{Model Setup and Calibration}

In the model, firms differ in productivity and face idiosyncratic shocks that may trigger financial distress. In default, they may be liquidated or reorganized. Reorganization allows firms to continue operating, but incurs a fixed cost, which introduces endogenous economies of scale in bankruptcy: large firms typically reorganize, while small firms often liquidate. Lenders anticipate these outcomes when pricing loans. Asset-based loans are secured against physical assets and less exposed to liquidation losses. In contrast, cash flow-based loans depend predominantly on expected future earnings, thus highly sensitive to liquidation probability. Consequently, credit spreads on CF-based loans rise sharply with firms' ex-ante likelihood of liquidation, endogenously linking reorganization costs to the pricing and composition of corporate debt. Elevated liquidation risk prevents productive but asset-poor firms from credibly pledging future earnings as collateral, sustaining credit misallocation and discouraging entry. The combination of these factors depresses aggregate productivity.

The model is disciplined using a broad dataset of U.S. non-financial firms between 2010 and 2024. The calibrated model replicates key facts about debt financing strategies and default resolution decisions for small and large U.S. firms. It also matches the differences between small and large firms in debt to collateral ratios, financing costs, and liquidation probabilities, even without explicitly targeting these moments - see Table \ref{tab:untargeted}.

\begin{table}[h!]
    \centering
    \begin{tabular}{lcccc}
    \toprule
     &  \multicolumn{2}{c}{\textbf{Data}} & \multicolumn{2}{c}{\textbf{Model}} \\ 
     \textbf{SS Value}  & Small Firms & Large Firms & Small Firms & Large Firms  \\ 
    \midrule
    Debt to Collateral & 0.44 & 0.59 & 0.45 & 0.57 \\ 
    Interest Rate (\%) & 5.89 & 4.63 & 5.63 & 4.89 \\ 
    Liquidation Probability & 0.84 & 0.27 & 0.85 & 0.14 \\ 
    CF Reliance & 0.51 & 0.53 & 0.24 & 0.97 \\ 
    CF-based to total debt & 0.34 & 0.83 & 0.30 & 0.87 \\ 
    \bottomrule
    \end{tabular}
    \caption{Untargeted data moments and their model counterparts}
    \label{tab:untargeted}
\end{table}

\section{Results}

To study the effects of  reorganization costs on the broader economy, I estimate the effects of an exogenous shock that cuts these costs by half. This could be interpreted as an exogenous improvement in reorganization technology or a policy reform that implements a more reorganization-friendly bankruptcy regime. Following such a shock, small firms, that are disproportionately burdened by reorganization costs, experience a drop in average liquidation probability from 86\% to 66\%, which fundamentally changes their optimal debt financing strategy. Their reliance on CF-based debt rises from 24\% to 43\%, allowing them to rely on external finance more intensively. Their debt-to-collateral ratio increases by 3 percentage points (from 0.45 to 0.48), and the average interest rate declines by 0.4 percentage points (from 5.63 to 5.23). Large firms experience qualitatively similar, but smaller changes, since for these firms liquidation risk is not a major impediment to borrowing. As a result, the financing gap between small and large firms narrows considerably.

The reduction in reorganization costs raises average productivity by 1.75\%, allowing firms to increase wages by a similar proportion. I decompose this productivity improvement into firm mass, composition, and capital deepening components. The model suggests that credit frictions reduce firm mass mainly by forcing low- and mid-productivity firms out of the market. As a result, the most productive firms that are largely unhindered by financial frictions capture an excessive share of the market. Lowering reorganization costs represents a meaningful shift toward the frictionless case, substantially increasing firm mass by allowing more middle-productivity firms to operate. Hence, the productivity gains are primarily driven by the increase in firm mass, while capital deepening contributes relatively little, and the composition effect (determined by the productivity of operating firms) partly offsets the gains.

\begin{figure}[H]
  \centering
  \textbf{ Firm  and Capital distribution Across Productivity States \vspace{2mm}} \\  % This acts as a title
  \includegraphics[width=1\textwidth]{C:/Users/szjud/OneDrive/Asztali gép/EBCs/CFL-git/Latex codes/Plots/Combined4runs2.png}
  \caption{\small The effects of lower reorganization costs on capital allocation and firm distribution across productivity states. Blue line: baseline allocation, black line: no-credit-frictions benchmark, red dashed line: allocation with 50\% reduction of reorganization costs; green dashed line: 90\% reduction.}
  \label{fig:firmalloc}
\end{figure}


Finally, I study the effects of the Small Business Reorganization Act (SBRA), introduced in the U.S. in order to make reorganization more accessible for small firms. Using the empirical estimates of Hotchkiss, Iverson, and Zheng (2024), I calibrate a targeted reduction in fixed reorganization costs to lower small firms' liquidation risk by about 12 percentage points. The model predicts that this reform raises aggregate productivity by 0.81\%. Extending the same reduction in reorganization costs to all firms increases the aggregate effect only marginally, indicating that size-based reforms are an efficient policy design. On the other hand, by weakening creditor control during renegotiation, the SBRA may induce lenders to raise spreads on cash flow–based debt, thereby partially offsetting the benefits of lower reorganization costs.

\section{Conclusion}

I study the macroeconomic implications of large reorganization costs in a general equilibrium model. By endogenously linking bankruptcy costs, debt composition, and credit pricing, the model captures both firm-level financing patterns and their macroeconomic consequences. Taken together, the results show that reorganization costs play a quantitatively important role in shaping credit allocation and aggregate productivity. Reorganization-friendly bankruptcy regimes that lower fixed costs of restructuring can substantially reduce misallocation and enhance productivity by improving small firms' access to external finance.



\section*{References}
	
\noindent
\hangindent=0.7cm
[1] Antill, S., \& Grenadier, S. R. (2019). Optimal capital structure and bankruptcy choice: Dynamic bargaining versus liquidation. \textit{Journal of Financial Economics}, 133(1), 198--224.
\vspace{3mm}

\noindent
\hangindent=0.7cm
[2] Corbae, D., \& D’Erasmo, P. (2021). Reorganization or liquidation: Bankruptcy choice and firm dynamics. \textit{The Review of Economic Studies}, 88(5), 2239--2274.
\vspace{3mm}

\noindent
\hangindent=0.7cm
[3] Drechsel, T. (2023). Earnings-based borrowing constraints and macroeconomic fluctuations. \textit{American Economic Journal: Macroeconomics}, 15(2), 1--34.
\vspace{3mm}

\noindent
\hangindent=0.7cm
[4] Hotchkiss, E., Iverson, B., \& Zheng, X. (2024). Can small businesses survive Chapter 11? \textit{Available at SSRN}.
\vspace{3mm}

\noindent
\hangindent=0.7cm
[5] Khan, A., Senga, T., \& Thomas, J. K. (2017). Default risk and aggregate fluctuations in an economy with production heterogeneity.
\vspace{3mm}

\noindent
\hangindent=0.7cm
[6] Lian, C., \& Ma, Y. (2021). Anatomy of corporate borrowing constraints. \textit{The Quarterly Journal of Economics}, 136(1), 229--291.
\vspace{3mm}

\noindent
\hangindent=0.7cm
[7] Tamayo, C. E. (2017). Bankruptcy choice with endogenous financial constraints. \textit{Review of Economic Dynamics}, 26, 225--242.
\vspace{3mm}



\end{document}